% chapters/14/section-04.tex
\section{Mengganti Greybox dengan Aset}
\label{sec:ganti-aset}

\subsection{Urutan Penggantian}

Penggantian balok abu-abu dengan aset sebaiknya dilakukan menurut urutan yang
tetap, agar tidak ada yang terlewat.

\begin{langkah}
  \item \textbf{Pastikan mekaniknya sudah selesai.} Seluruh ukuran pada
        Bagian~13.1 telah ditetapkan, dan levelnya sudah dapat dimainkan
        dengan lancar.
  \item \textbf{Impor seluruh asetnya lebih dahulu}, lalu atur
        \istilah{importer setting}-nya satu per satu. Mengatur pengaturan
        sekaligus pada awal pekerjaan lebih cepat daripada mengaturnya
        sambil mengganti benda.
  \item \textbf{Ganti benda satu per satu}, mulai dari yang paling terlihat
        pemain: pemain, lalu rintangan, lalu latar. Setelah setiap
        penggantian, jalankan scene-nya dan pastikan mekaniknya tidak berubah.
  \item \textbf{Periksa bentuk batas bendanya.} Mengganti gambar benda tidak
        mengubah \istilah{Collider}-nya secara otomatis. Apabila ukuran
        gambarnya berbeda dari balok sebelumnya, bentuk batasnya perlu
        disesuaikan.
\end{langkah}

Langkah keempat itu paling sering terlewat, dan gejalanya khas: benda yang
sudah bergambar rapi tetap bertabrakan seolah-olah masih berupa balok
sebelumnya.

\subsection{Membandingkan Sebelum dan Sesudah}

Perbandingan sebelum dan sesudah penggantian aset merupakan bagian dari luaran
praktikum pada RPS. Perbandingan itu berguna untuk dua hal.

Pertama, untuk memastikan bahwa mekaniknya \textbf{tidak berubah}. Apabila
pemain tiba-tiba tidak dapat melompati celah yang sebelumnya terjangkau,
kemungkinan besar bentuk batas bendanya berubah bersama gambarnya.

Kedua, untuk memeriksa keputusan perancangan Anda. Sering kali aset yang
sesungguhnya menuntut penyesuaian pada levelnya: celah yang tadinya terasa
wajar menjadi terasa terlalu sempit, atau warna latar yang tadinya netral
menjadi bertabrakan dengan warna pemain. Temuan semacam itu wajar, dan justru
merupakan alasan mengapa level disusun lebih dahulu secara
\istilah{greyboxing}.

\begin{tipsbox}
Simpan tangkapan layar sebelum dan sesudah penggantian aset, dan letakkan
keduanya berdampingan pada laporan praktikum. Kebiasaan mendokumentasikan
perubahan semacam itu berguna juga pada proyek akhir nanti, sebab laporan
pengembangan yang memuat bukti perubahan lebih meyakinkan daripada yang hanya
menceritakannya.
\end{tipsbox}

\begin{peringatanbox}
Apabila Anda memakai aset buatan orang lain, catatlah sumbernya beserta
lisensinya sejak awal, ketika asetnya baru diimpor. Mencari kembali asal-usul
ratusan berkas pada akhir semester jauh lebih melelahkan daripada mencatatnya
satu per satu ketika diimpor. Cara memilih aset beserta kewajiban atribusinya
dibahas pada Bagian~\ref{sec:hak-aset}.
\end{peringatanbox}
